% SPDX-License-Identifier: GPL-3.0-or-later
% Copyright (C) 2026 Alessandro Greco
%% Русское альтернативное описание EPUB.
\altfig{На схеме цепь клиент — сторожевой — промежуточный — выходной ретранслятор проходит три канала с разными ненулевыми локальными CircID: 0x8000A1C3 на первом, 0x0000712F на втором и 0x93B40E07 на третьем. Таблица сторожевого узла связывает первый идентификатор со вторым, а таблица промежуточного — второй с третьим. На каналах с протоколом v4+ CircID занимает 32 бита; старший бит зависит от того, кто инициировал соединение канала, а не от номера перехода цепи. Сравнение одних CircID не позволяет напрямую сопоставить каналы, но не исключает корреляции по другим признакам трафика.}
% CircID имеет локальную область действия внутри канала.
\begin{tikzpicture}[node distance=2.3cm]
  \node[client] (cl) {Клиент};
  \node[relay,right=of cl] (g) {Сторожевой};
  \node[relay,right=of g] (m) {Промежуточный};
  \node[relay,right=of m] (x) {Выходной};
  \draw[thick,torgrigio] (cl) -- node[above,font=\sffamily\scriptsize]{канал 1} node[below,font=\ttfamily\scriptsize,text=torviola]{0x8000A1C3} (g);
  \draw[thick,torgrigio] (g) -- node[above,font=\sffamily\scriptsize]{канал 2} node[below,font=\ttfamily\scriptsize,text=torviola]{0x0000712F} (m);
  \draw[thick,torgrigio] (m) -- node[above,font=\sffamily\scriptsize]{канал 3} node[below,font=\ttfamily\scriptsize,text=torviola]{0x93B40E07} (x);
  \node[below=0.9cm of g,draw=torgrigio,fill=white,font=\ttfamily\scriptsize,align=left] (tg)
    {(кан.1, 0x8000A1C3)\\\hspace*{1em}$\leftrightarrow$ (кан.2, 0x0000712F)};
  \node[below=0.9cm of m,draw=torgrigio,fill=white,font=\ttfamily\scriptsize,align=left] (tm)
    {(кан.2, 0x0000712F)\\\hspace*{1em}$\leftrightarrow$ (кан.3, 0x93B40E07)};
  \draw[dotted] (g) -- (tg); \draw[dotted] (m) -- (tm);
\end{tikzpicture}
